% Sección 11.2: Vectores

\section*{SECCIÓN 11.2 \quad Vectores}

Los científicos utilizan el término vector para indicar una cantidad (tal como velocidad, fuerza o aceleración) que posee tanto magnitud como dirección.

\vspace{0.5cm}
\noindent\textit{En los ejercicios 1-6 encuentre un vector $\mathbf{a}$ con representación dada por el segmento lineal dirigido $\overrightarrow{AB}$. Dibuje $\overrightarrow{AB}$ y la representación equivalente que empieza en el origen.}

\begin{enumerate}
    \subimport{ejercicio_01/}{ejercicio.tex}
\subimport{ejercicio_02/}{ejercicio.tex}
\subimport{ejercicio_03/}{ejercicio.tex}
\subimport{ejercicio_04/}{ejercicio.tex}
\subimport{ejercicio_05/}{ejercicio.tex}
\subimport{ejercicio_06/}{ejercicio.tex}
\end{enumerate}

\vspace{0.3cm}
\noindent\textit{En los ejercicios 7-10 encuentre la suma de los vectores dados y haga una ilustración geométrica.}

\begin{enumerate}
    \setcounter{enumi}{6}
    \subimport{ejercicio_07/}{ejercicio.tex}
\subimport{ejercicio_08/}{ejercicio.tex}
\subimport{ejercicio_09/}{ejercicio.tex}
\subimport{ejercicio_10/}{ejercicio.tex}
\end{enumerate}

\vspace{0.3cm}
\noindent\textit{En los ejercicios 11-18 encuentre $|\mathbf{a}|$, $\mathbf{a} + \mathbf{b}$, $\mathbf{a} - \mathbf{b}$, $2\mathbf{a}$ y $3\mathbf{a} + 4\mathbf{b}$.}

\begin{enumerate}
    \setcounter{enumi}{10}
    \subimport{ejercicio_11/}{ejercicio.tex}
\subimport{ejercicio_12/}{ejercicio.tex}
\subimport{ejercicio_13/}{ejercicio.tex}
\subimport{ejercicio_14/}{ejercicio.tex}
\subimport{ejercicio_15/}{ejercicio.tex}
\subimport{ejercicio_16/}{ejercicio.tex}
\subimport{ejercicio_17/}{ejercicio.tex}
\subimport{ejercicio_18/}{ejercicio.tex}
\end{enumerate}

\vspace{0.3cm}
\noindent\textit{En los ejercicios 19-24 encuentre un vector unitario que tenga la misma dirección del vector dado.}

\begin{enumerate}
    \setcounter{enumi}{18}
    \subimport{ejercicio_19/}{ejercicio.tex}
\subimport{ejercicio_20/}{ejercicio.tex}
\subimport{ejercicio_21/}{ejercicio.tex}
\subimport{ejercicio_22/}{ejercicio.tex}
\subimport{ejercicio_23/}{ejercicio.tex}
\subimport{ejercicio_24/}{ejercicio.tex}
\end{enumerate}

\vspace{0.3cm}
\noindent\textit{En los ejercicios 25 y 26 exprese $\mathbf{i}$ y $\mathbf{j}$ en términos de $\mathbf{a}$ y $\mathbf{b}$.}

\begin{enumerate}
    \setcounter{enumi}{24}
    \subimport{ejercicio_25/}{ejercicio.tex}
\subimport{ejercicio_26/}{ejercicio.tex}
\end{enumerate}

\vspace{0.3cm}
\begin{enumerate}
    \setcounter{enumi}{26}
    \subimport{ejercicio_27/}{ejercicio.tex}
\subimport{ejercicio_28/}{ejercicio.tex}
\subimport{ejercicio_29/}{ejercicio.tex}
\subimport{ejercicio_30/}{ejercicio.tex}
\subimport{ejercicio_31/}{ejercicio.tex}
\subimport{ejercicio_32/}{ejercicio.tex}
\subimport{ejercicio_33/}{ejercicio.tex}
\subimport{ejercicio_34/}{ejercicio.tex}
\subimport{ejercicio_35/}{ejercicio.tex}
\end{enumerate}
